\documentclass[11pt]{article}
\usepackage{icmlko}
\begin{document}

\title{어제 세운 조항을 오늘 어겼다 --- 그래서 조항을 사람에게서 코드로 옮겼다}
\author{월드모델 실험실}
\date{노트 560 · 561 · 2026년 8월 5일}
\maketitle

\begin{abstract}
노트 $559$ 의 도메인별 표에서 \textbf{도서만 $+0.0459$} 로 올랐다. 도메인별
\texttt{zero\_is\_data} 를 사전 등록하고 도서만 바꿨다.

\begin{center}
\begin{tabular}{lrr}
\hline
 & 도서 $\rho$ 변화 & 판 \\
\hline
전 도메인 False(노트 $559$) & $\mathbf{+0.0459}$ & $-0.0034$ \\
\textbf{도서만 False}(여기) & $\mathbf{-0.0820}$ ($0/12$) & $-0.0033$ \\
도서 검색 통째 가림(대조) & $-0.1493$ ($0/12$) & $-0.0075$ \\
\hline
\end{tabular}
\end{center}

\textbf{부호가 뒤집힌다.} 도서의 $+0.0459$ 는 도서 자신의 변경이 아니라
\emph{이웃 효과}였다.

\begin{quote}
\textbf{그리고 이것은 노트 $551$ 이 세운 조항이다} --- ``\emph{전역으로 끄면
오른다}''를 ``\emph{거기서 가리면 오른다}''로 읽지 않는다. 전역 제거의
도메인별 차에는 이웃 효과가 섞여 있다. \textbf{그 조항을 쓴 다음 날 내가
어겼다.}
\end{quote}
\end{abstract}

\begin{figure}[t]
\centering
\includegraphics[width=\textwidth]{figs/f1.pdf}
\caption{\textbf{가} --- 같은 함정, 두 번. \textbf{나} --- 둘 다 음수.
\textbf{다} --- 두 경로.}
\end{figure}

\section{경로가 둘이다}

한 도메인의 예보가 바뀌는 길은 둘이다.

\begin{enumerate}
\item 자기 축이 없어져서 --- \textbf{그 도메인만 바꿔도 난다}
\item 남들이 그 축으로 배운 것이 없어져서 --- \textbf{전역으로 껐을 때만 난다}
\end{enumerate}

전역 제거의 도메인별 차는 ①$+$②다. 노트 $551$ 의 팝업이 ②의 순수한
사례였고($+0.0374 \to +0.0018$, 팝업은 애초에 검색 축이 없다), 도서는
\textbf{①과 ②의 부호가 반대}인 사례다 --- ②가 $+$, ①이 $-$, 합이 $+$.

\section{대조 팔이 가설도 죽였다}

사전 등록에 \textbf{갈라 볼 것}을 적어 뒀다 --- 도서가 오르는 것이 (가)
관측-$0$ 이 해로워서인지 (나) 그냥 축이 줄어서인지. (나)라면 검색을 통째로
가려도 올라야 한다.

\begin{center}
\begin{tabular}{lr}
\hline
 & 도서 $\rho$ 변화 \\
\hline
B --- 관측-$0$ 만 뺌 & $-0.0820$ \\
C --- 검색 통째 가림 & $-0.1493$ \\
\hline
\end{tabular}
\end{center}

C 는 노트 $548$ 이 정정 전에 잰 $-0.1654$ 와 맞는다. \textbf{둘 다 음수다.}

\begin{quote}
\textbf{가설 (가)도 반증됐다.} 관측-$0$ 만 빼도 $-0.0820$ 이므로
\textbf{관측-$0$ 도 정보}다 --- ``긁었는데 검색량이 문턱 아래''가 곧
``무명하다''이고, 도서에서 그것이 세다. 노트 $548$ 이 잰 결측
표시자$\leftrightarrow$라벨 $+0.1808$(도서)이 같은 이야기였다.
\end{quote}

\section{사람이 기억하는 조항은 하루를 못 갔다}

\begin{center}
\begin{tabular}{lll}
\hline
 & 노트 $551$(팝업) & 노트 $560$(도서) \\
\hline
전역 차 & $+0.0374$ & $+0.0459$ \\
그 도메인만 & $+0.0018$ (문다) & $\mathbf{-0.0820}$ \\
\hline
\end{tabular}
\end{center}

노트 $551$ 이 조항을 쓰고 논문 $432$ 까지 냈다. 그 다음 노트에서 같은 종류의
표(노트 $559$ 의 도메인별 차)를 보고 \textbf{같은 방식으로 후보를 세웠다.}

\begin{quote}
\textbf{그래서 조항을 코드로 옮겼다}(노트 $561$).
\texttt{sideaudit.domain\_deltas} 가 도메인별 차 표를 낼 때 경고를 함께
낸다 --- ``\emph{이 표로 가림 후보를 세우지 않는다}''. 이 함수를 거치지 않고
그런 표를 인쇄하지 않는다.
\end{quote}

\section*{같은 커밋에서 고친 것 하나 더}

검사 ③ 이 자료만 보고 순위를 내서 \textbf{노트 $553$ 이 채택한 만화
\texttt{trend} 가 여전히 $2$ 위}로 남아 있었다. \texttt{\_handled()} 가
손잡이 셋(\texttt{DOMDROP} $\cdot$ \texttt{BLOCK} $\cdot$
\texttt{TREND\_DROP}/\texttt{WIKI\_DROP})을 한꺼번에 읽어 ``\textbf{이미
처리됨}''을 표시하고, ``\textbf{아직 안 고친 것 $N$ 개}''를 따로 센다.

\begin{quote}
지금 안 고친 것은 다섯이고 제일 큰 것은 \textbf{웹툰 trend}(곱 $0.043$)다.
그런데 노트 $558$ 의 표가 웹툰은 \textbf{값으로 번다}($+0.1809$)고 말한다 ---
곱은 화면을 좁히는 데 쓰고 판정은 그 표로 한다.
\end{quote}

\section*{조항}

\begin{quote}
\textbf{조항이 두 번 깨지면 사람에게 두지 않는다.} 노트 $551$ 이 세우고
노트 $560$ 이 어겼다 --- 사이에 하루도 없었다. \textbf{표를 만드는 함수가
경고를 같이 내게 한다.}
\end{quote}

\section*{남은 것}

\begin{itemize}
\item \textbf{도서의 관측-$0$ 이 왜 그렇게 센가} --- 통째 가림 $-0.1493$ 중
관측-$0$ 몫이 $-0.0820$ 으로 \textbf{절반이 넘는다}. 도서는 관측 축이 여섯
뿐이라(노트 $548$) 검색이 대신할 것이 없다는 읽기와 맞는데, 직접 안 봤다.
\item \textbf{이웃 효과를 잡는 손잡이가 없다} --- ①은 \texttt{DOMDROP} 으로
잡지만 ②는 못 잡는다. 도메인별 모형이나 가중이 필요한데 그것은 다른 층이다.
\item \textbf{노트 $559 \cdot 560$ 이 연달아 내 사전 등록의 결함}이었다 ---
$559$ 는 전제(``이것은 정정이다''), $560$ 은 후보 선정. \textbf{둘 다 측정이
아니라 읽기에서 틀렸다.}
\end{itemize}

\end{document}
